\documentclass[11pt]{amsart}
\usepackage[a4paper,margin=2.5cm]{geometry}
\usepackage{amssymb,amsthm,booktabs,xcolor}
\usepackage[hidelinks]{hyperref}

% placeholders: what a section still needs (shown in blue in the PDF; delete when written)
\newcommand{\plan}[1]{\par\noindent{\color{blue}\small[\emph{To write:} #1]}\par}
% points to cover, as a blue list (delete each list once its section is written)
\newenvironment{points}[1]{\par\noindent{\color{blue}\small\emph{Points to cover: #1}}\color{blue}\small\begin{itemize}\setlength\itemsep{1pt}}{\end{itemize}\normalcolor\normalsize}

\newtheorem{theorem}{Theorem}
\newtheorem{proposition}[theorem]{Proposition}
\newtheorem{lemma}[theorem]{Lemma}
\theoremstyle{remark}
\newtheorem{remark}[theorem]{Remark}

\title{Searching for six coprime representations as a sum of two cubes}
\author{David Curran}
\date{\today}
\address{University College Dublin, Ireland} 
\email{david.curran1@ucd.ie}
\begin{document}

\begin{abstract}
The sequence of least positive integers with $k$ representations as $x^3 + y^3$ with $\gcd(x, y) = 1$, where $x$
and $y$ may be negative, is known up to the fifth term. This number $506433677359393$ was found by Bernstein in 1999, and the sixth is unknown. It is sequence A047697 in the OEIS, and we report two searches for its sixth term. The first search is exhaustive: it uses GPUs to search every $n \le 2.5\cdot10^{23}$. The search proves
A047697(6) $> 2.5\cdot10^{23}$, a millionfold beyond Bernstein's bound, and finds exactly twelve coprime five-ways in
that range. Since a cubefree number has only coprime representations, it also gives a lower bound for the cubefree taxicab numbers (OEIS A080642):
A080642(5) $> 2.5\cdot10^{23}$. The second search uses the geometry of the elliptic curve $x^3 + y^3 = n$: in every
example we tested ($1{,}378$ numbers), the integer points are related only by lines (three collinear points) and hubs
(two chords meeting at the same rational point), so we enumerate lines and hubs directly. It is not exhaustive but
reaches far larger $n$, and it has found 70 further coprime five-ways, up to $3.8\cdot10^{35}$. We prove that no line
carries three points with positive coordinates, so a cubefree number with many positive representations must come
from hubs or high rank; a search restricted to positive hubs finds eight cubefree numbers with four positive
representations. We found no six-way, and no cubefree five-way. We release the catalogue of $8{,}788$ coprime
four-ways and 82 five-ways, complete for $n \le 10^{22}$ and $10^{23} < n \le 2.5\cdot10^{23}$, together with the
code.
\end{abstract}

\maketitle

\section{Introduction}
\begin{points}{the problem}
\item Hardy and Ramanujan, 1729; the taxicab numbers Ta$(k)$ (OEIS A011541) and cabtaxi numbers Ca$(k)$ (A047696, negative cubes allowed).
\item Ta$(k)$ exists for every $k$ (Hardy and Wright, Theorem 412), but only by scaling rational points by cubes; the numbers it gives are far from coprime.
\item The coprime version A047697 and the cubefree version A080642 forbid scaling, so a $k$-way needs $k$ genuinely distinct integer points. Whether they exist for every $k$ is open; Silverman \cite{silverman1993} links it to the rank of $x^3+y^3=n$ being unbounded; Stewart \cite{stewart2008,stewart2009} on cubic Thue equations with many solutions.
\item Known terms: A047697 = 1, 91, 3367, 16776487 (Rathbun), 506433677359393 (Bernstein); A080642 = 2, 1729, 15170835645, 1801049058342701083 (Hollerbach's list).
\item Relatives in the OEIS: A023050, A293645--A293651 (Fay, 2017).
\end{points}
\begin{points}{best current results}
\item Coprime: Bernstein \cite{bernstein2001} found A047697(5) and searched to $2.5\cdot10^{17}$; no later search known.
\item Cubefree: A080642(4) from Hollerbach's search; no published bound for A080642(5).
\item Taxicab: Wilson \cite{wilson1999} Ta(5); Calude, Calude and Dinneen \cite{caludedinneen2003} Ta(6) with high probability; Hollerbach \cite{hollerbach2008} announced an exhaustive confirmation; Eisermann \cite{eisermann2008} Ta(6) $>5\cdot10^{20}$.
\item Upper bounds for Ta$(k)$ and Ca$(k)$: Boyer \cite{boyer2008}, Su \cite{su2016} (context only; our own bounds are a separate paper).
\item High rank in this family: Elkies and Rogers \cite{elkiesrogers2004}.
\end{points}
\begin{points}{websites}
\item OEIS entries \cite{oeisA047697,oeisA080642}; Boyer's taxicab pages (current bound tables); Butler's pages \cite{butler} (primitive lists complete to $10^{23}$). Cite papers for results, websites for data.
\end{points}
\begin{points}{two families of algorithm}
\item Exhaustive enumeration: list every $a^3\pm b^3\le B$ and find repeats; all methods cost about $B^{2/3}$ (the number of pairs) and differ in how repeats are found: a heap in increasing order (Wilson, Bernstein; memory $B^{1/3}$), bimonotone enumeration (Eisermann), bands plus a hash table (Butler).
\item Construction from rational points: take $x^3+y^3=M$ with many rational points and clear denominators with a cube (Boyer, Su, Butler above $10^{23}$). Gives the Ta/Ca upper bounds, but never a coprime or cubefree number, since the scaling factor divides every representation.
\item So the coprime problem needs brute force or a different structure: this is the motivation for our two searches.
\end{points}
\begin{points}{what this paper does}
\item Exhaustive GPU search to $2.5\cdot10^{23}$: A047697(6) $>2.5\cdot10^{23}$, hence A080642(5) $>2.5\cdot10^{23}$; exactly 12 five-ways in that range.
\item The geometry: only lines and hubs relate the integer points; theorems that make lines searchable; convexity (no line holds three positive points).
\item A targeted search by lines and hubs: 70 further five-ways up to $3.8\cdot10^{35}$; eight all-positive cubefree four-ways; no six-way.
\item Catalogue of 8,788 four-ways and 82 five-ways; OEIS extensions (A023050, A293648, A293649, A293651 with 14 missing terms).
\item Outline of the sections.
\end{points}

\section{The exhaustive search}
\begin{points}{what is counted}
\item Every coprime representation of $n\le B$, written once: sums $a^3+b^3$ with $1\le a\le b$, and differences $(c+d)^3-c^3 = d^3+3dc(c+d)$ with $c,d\ge1$.
\item About $1.33\,B^{2/3}$ pairs, sums about a third of them; $6/\pi^2\approx 60.8\%$ coprime (Baker \cite{baker2007}, Liu \cite{liu2015} for the count).
\end{points}
\begin{points}{the algorithm}
\item Same family as Butler's: cut $(0,B]$ into segments, generate every pair whose value falls in a segment, find repeats inside the segment. Butler hashes; we sort on a GPU.
\item Segments equal in $n^{2/3}$ (equal numbers of pairs); a segment that is too large is split; segments are independent, so $k$ GPUs take equal shares of the work.
\item On the GPU, per segment: count (one thread per stream, exact cube roots or an exact quadratic bound), prefix-sum, write (additions only; values stored as 64-bit offsets $n-\mathrm{lo}$, so 128-bit values sort as 64-bit keys; binary gcd removes non-coprime pairs), radix sort, mark runs of length $\ge k_{\min}$.
\item On the CPU, every reported $n$ re-derived independently in exact 128-bit arithmetic; the count must equal the GPU's multiplicity.
\item Number theory only for checking, never for pruning (THEOREMS 5; the divisor test, cf.\ Broughan \cite{broughan2003}), so ``no six-way up to $B$'' remains a proof.
\end{points}
\begin{points}{validation}
\item Brute force to $10^7$; known terms; splitting a range gives the same pairs; a window above $2^{64}$; agreement with the CPU program at $10^{12}$ and $10^{15}$; self-test on every GPU before every run; every $k\ge4$ number recounted from its factorisation; 0--$10^{22}$ run twice with identical output.
\end{points}
\begin{points}{runs, cost and results}
\item Table of ranges, hardware, pairs and time (about 140 GPU-hours on RTX 5090s, plus 7of9's RTX 2080 Ti runs); sort is 60\% of the time.
\item Results: A047697(6) $>2.5\cdot10^{23}$; the 12 five-ways by decade (none in $10^{22}$--$10^{23}$); complete four-way lists.
\item Comparison: Bernstein $2.5\cdot10^{17}$; Eisermann's taxicab enumeration to $5\cdot10^{20}$ (about 650 CPU-days); Butler complete to $10^{23}$ (positive, primitive).
\item Limits and next steps: the $10^{24}$ guard; bucket counting instead of a full sort (Butler's hash idea on the GPU), about $2\times$.
\end{points}

\section{The geometry of \texorpdfstring{$x^3 + y^3 = n$}{x^3 + y^3 = n}}
\plan{the curve and its flex at infinity; collinear points sum to zero, $-P$ is the swap; lines, bowties, hubs,
D-lines; the census of relations (only lines and hubs); structure against rank. Reference notes:
\texttt{writeup/02\_curve\_search.md} \S1.}

\section{Lines that can carry three coprime points}
\plan{Theorem 1 (direction rule), Theorem 2 ($E \mid n$), Theorem 4 (complete enumeration via the conic), Theorem 7
(D-lines), Theorem 8 (types at 2 and 3; no complete quadrilateral), convexity. Proofs from
\texttt{experiments/mining/THEOREMS.md}.}

\section{The curve search}
\begin{points}{the idea}
\item Most multi-ways are tied together by lines (three collinear integer points) and hubs; so instead of enumerating pairs, enumerate lines and hubs, and count every number they produce exactly.
\item Not exhaustive: it finds numbers, it cannot rule them out. Reaches $n$ far beyond the GPU search.
\end{points}
\begin{points}{lines}
\item Directions $(a,b)$ of lines $ax+by=c$ that can carry three coprime points: Theorem 1 ($ab>0$, a norm from $\mathbb{Q}(\sqrt{-3})$, $4\mid ab$); cost $b^3-a^3$; all allowed directions in increasing cost.
\item For each direction, the collinear triples are the rational points of the conic $X^2+XY+Y^2=3ab^3$ (Theorem 4): sweep the slopes $r/s$ through one point found by \texttt{qfsolve}, up to a height; integer arithmetic throughout.
\item Height per direction from $n\approx \mathrm{cost}^2h^6$, stopping where $n$ passes the cap (keeps 94\% of candidates in 36\% of the time).
\end{points}
\begin{points}{filter and exact counting}
\item Keep $n$ with at least 7 split primes ($p\equiv1 \bmod 3$) below 1,000 and two of $\{7,13,19,37\}$; why (five-ways are rich in split primes; the free-point rate doubles per extra prime).
\item Count every kept $n$ exactly from its factorisation: for each divisor $s\le(4n)^{1/3}$ test whether $(4n/s-s^2)/3$ is a square (about $10^4$ numbers per second per core). Any extra representation is found whatever its shape.
\end{points}
\begin{points}{D-lines and hubs}
\item A D-line: two integer points and a third of denominator $D$; scaled by $D$ it is an ordinary line on $D^3n$ (Theorem 7); two D-lines through the same point form a hub.
\item Only one chord of a hub has to be enumerated; exact counting finds the other.
\item Positive hubs: the slopes whose line has two positive points form three arcs; enumerating only those is about $5\times$ cheaper and gives the all-positive hubs needed for the cubefree problem.
\end{points}
\begin{points}{runs and saturation}
\item Machines and runs (Mac, 7of9, worf, rented GPUs for the first line runs); candidates counted.
\item Saturation below $10^{30}$: every direction of cost $\le1.2\cdot10^{12}$ searched; most new five-way finds became rediscoveries.
\item Recall against the exhaustive data (50--70\% of four-ways below $10^{16}$, falling with $n$; higher for five-ways).
\end{points}

\section{Results}
\plan{the 82 five-ways (table in the appendix); four-way counts; recall of the curve search against exhaustive data;
shapes.}

\section{Cubefree taxicab numbers}
\plan{cubefree implies coprime; the bound A080642(5) $> 2.5\cdot10^{23}$; convexity rules out lines; the eight
all-positive four-ways; why a five-way is out of reach of every route we tried.}

\section{Why no six-way?}
\plan{the shapes a six-way could take; the free-point rate and the estimate of one six-way per 300--1,500 five-ways;
what a successful search would need.}

\section{Data and code}
\plan{repository, catalogue format, how to reproduce each claim.}

\bibliographystyle{amsplain}
\bibliography{refs}

\end{document}
